\documentclass[12pt,a4paper]{article}

% ── Packages ──────────────────────────────────────────────────────────────────
\usepackage[utf8]{inputenc}
\usepackage[T1]{fontenc}
\usepackage[margin=2.5cm]{geometry}
\usepackage{hyperref}
\usepackage{amsmath}
\usepackage{amssymb}
\usepackage{graphicx}
\usepackage{booktabs}
\usepackage{xcolor}
\usepackage{listings}
\usepackage{enumitem}
\usepackage{titlesec}
\usepackage{fancyhdr}
\usepackage{parskip}
\usepackage{microtype}
\usepackage{caption}
\usepackage{float}
\usepackage{mdframed}

% ── Colour palette ────────────────────────────────────────────────────────────
\definecolor{javakw}{RGB}{127,0,85}      % keywords  (Eclipse dark purple)
\definecolor{javastr}{RGB}{42,0,255}     % strings   (blue)
\definecolor{javacomment}{RGB}{63,127,95}% comments  (green)
\definecolor{codebg}{RGB}{248,248,248}
\definecolor{coderule}{RGB}{200,200,200}
\definecolor{accent}{RGB}{0,90,160}

% ── Java listings style ───────────────────────────────────────────────────────
\lstdefinestyle{java}{
  language=Java,
  backgroundcolor=\color{codebg},
  basicstyle=\ttfamily\small,
  keywordstyle=\color{javakw}\bfseries,
  stringstyle=\color{javastr},
  commentstyle=\color{javacomment}\itshape,
  numberstyle=\tiny\color{gray},
  numbers=left,
  numbersep=8pt,
  frame=single,
  rulecolor=\color{coderule},
  breaklines=true,
  breakatwhitespace=true,
  showstringspaces=false,
  tabsize=4,
  captionpos=b,
  morekeywords={var,record,sealed},
}
\lstdefinestyle{plain}{
  language={},
  backgroundcolor=\color{codebg},
  basicstyle=\ttfamily\small,
  keywordstyle={},
  stringstyle={},
  commentstyle={},
  numberstyle=\tiny\color{gray},
  numbers=left,
  numbersep=8pt,
  frame=single,
  rulecolor=\color{coderule},
  breaklines=true,
  showstringspaces=false,
  tabsize=4,
  captionpos=b,
}
\lstdefinestyle{bash}{
  language=bash,
  backgroundcolor=\color{codebg},
  basicstyle=\ttfamily\small,
  keywordstyle=\color{javakw}\bfseries,
  commentstyle=\color{javacomment}\itshape,
  numberstyle=\tiny\color{gray},
  numbers=left,
  numbersep=8pt,
  frame=single,
  rulecolor=\color{coderule},
  breaklines=true,
  showstringspaces=false,
  tabsize=2,
  captionpos=b,
}
\lstset{style=java}

% ── Header / footer ───────────────────────────────────────────────────────────
\pagestyle{fancy}
\fancyhf{}
\lhead{\textcolor{accent}{\small JIMP2 — Java Implementation}}
\rhead{\small\thepage}
\renewcommand{\headrulewidth}{0.4pt}

% ── Section formatting ────────────────────────────────────────────────────────
\titleformat{\section}{\large\bfseries\color{accent}}{{\thesection}}{1em}{}[\titlerule]
\titleformat{\subsection}{\normalsize\bfseries}{{\thesubsection}}{1em}{}

% ── Hyperref ──────────────────────────────────────────────────────────────────
\hypersetup{
  colorlinks=true,
  linkcolor=accent,
  urlcolor=accent,
  pdfauthor={Nghia Bao Phuc Ho},
  pdftitle={JIMP2 Java Implementation Documentation},
}

% ─────────────────────────────────────────────────────────────────────────────
\begin{document}

\begin{titlepage}
  \centering
  \vspace*{3cm}
  {\Huge\bfseries\color{accent} Implementation Documentation}\\[0.8em]
  {\Large JIMP2 — Java Edition}\\[2em]
  \rule{0.6\textwidth}{1pt}\\[2em]
  {\large Nghia Bao Phuc Ho}\\[0.5em]
  {\normalsize May 2026}\\[3em]
  \begin{mdframed}[backgroundcolor=codebg,linecolor=coderule]
    \centering
    \texttt{Graph Visualizer} · Java 17 · Swing GUI\\[0.3em]
    Algorithms: Fruchterman-Reingold \& Tutte's Embedding
  \end{mdframed}
  \vfill
  {\small
    This document describes the Java implementation of the JIMP2 graph
    visualisation project. It covers the data model, layout algorithms,
    I/O layer, and interactive Swing GUI. Where relevant, design choices
    are compared with the companion C implementation.}
\end{titlepage}

\tableofcontents
\newpage

% ─────────────────────────────────────────────────────────────────────────────
\section{Architecture Overview}
% (No direct counterpart in the C doc — Java's OOP warrants its own section)

Unlike the C implementation, which is structured as a flat collection of
\texttt{.c} and \texttt{.h} files linked into a single binary, the Java
project is divided into four packages with clearly separated
responsibilities:

\begin{center}
\begin{tabular}{lll}
\toprule
\textbf{Package} & \textbf{Classes} & \textbf{Responsibility} \\
\midrule
\texttt{model}     & \texttt{Node}, \texttt{Edge}, \texttt{Graph}          & Pure data, no logic \\
\texttt{io}        & \texttt{GraphReader}, \texttt{GraphWriter}             & File I/O \\
\texttt{algorithm} & \texttt{LayoutAlgorithm}, \texttt{Fr}, \texttt{Tutte} & Layout algorithms \\
\texttt{gui}       & \texttt{MainWindow}, \texttt{GraphPanel},             & Interactive GUI \\
                   & \texttt{ToolPanel}, \texttt{CoordinatePanel}          & \\
\bottomrule
\end{tabular}
\end{center}

\medskip
The application is single-threaded by design: all Swing events and
layout computations run on the Event Dispatch Thread (EDT) via
\texttt{SwingUtilities.invokeLater}.

% ─────────────────────────────────────────────────────────────────────────────
\section{Data Model}

\subsection{Node}

\texttt{Node} (in \texttt{model/Node.java}) is the Java equivalent of the
C \texttt{vertex\_t} struct. It stores a unique integer \texttt{id} and
mutable Cartesian coordinates \texttt{X}, \texttt{Y} that are updated
in-place by the layout algorithms.

\begin{lstlisting}[caption={Node class}]
public class Node {
    public int    id;
    public double X;
    public double Y;

    public Node(int id, double X, double Y) {
        this.id = id;  this.X = X;  this.Y = Y;
    }
}
\end{lstlisting}

\textbf{Design note.} All fields are \texttt{public} — the same philosophy
as C struct fields — because the algorithms and renderer need direct, fast
access without getter overhead. For a production codebase one would
encapsulate these, but the lab context favours simplicity.

\subsection{Edge}

\texttt{Edge} (in \texttt{model/Edge.java}) maps to \texttt{edge\_t} in the
C version. The \texttt{len} field carries the same semantic role as
\texttt{weight} in C: a multiplier applied to the attractive force in
Fruchterman-Reingold.

\begin{lstlisting}[caption={Edge class}]
public class Edge {
    public String name;
    public int    startNode;
    public int    endNode;
    public double len;          // edge weight (attractive-force multiplier)

    public Edge(String name, int st, int en, double len) {
        this.name = name;
        this.startNode = st;
        this.endNode   = en;
        this.len       = len;
    }
}
\end{lstlisting}

\subsection{Graph}

\texttt{Graph} (in \texttt{model/Graph.java}) is the top-level container.
It replaces the C \texttt{graph\_t} struct together with its explicit
\texttt{n\_vertices} and \texttt{n\_edges} counters; Java's
\texttt{ArrayList} tracks size automatically.

\begin{lstlisting}[caption={Graph class}]
public class Graph {
    public List<Node> nodes;
    public List<Edge> edges;

    public Graph() {
        this.nodes = new ArrayList<>();
        this.edges = new ArrayList<>();
    }
}
\end{lstlisting}

\textbf{Memory lifecycle.} Because Java uses garbage collection, there is
no manual \texttt{graph\_free} equivalent. The GC reclaims all
\texttt{Node}, \texttt{Edge}, and \texttt{Graph} objects once the
\texttt{GraphPanel} drops its reference to the old graph when a new file
is loaded.

% ─────────────────────────────────────────────────────────────────────────────
\section{Graph Loading — \texttt{GraphReader}}

\texttt{GraphReader.read(String filename)} in \texttt{io/GraphReader.java}
is the Java counterpart of \texttt{graph\_load} in C. It reads an
edge-list file and returns a fully populated \texttt{Graph} object, or
throws \texttt{FileNotFoundException} if the file cannot be opened.

\subsection{Algorithm}

The reader uses Java's \texttt{Scanner} for line-by-line parsing and a
\texttt{HashSet<Integer>} to track which node IDs have already been seen,
so each node is added to \texttt{graph.nodes} at most once.

\begin{lstlisting}[caption={GraphReader — core parsing loop}]
public Graph read(String filename) throws FileNotFoundException {
    Graph graph = new Graph();
    Scanner sc  = new Scanner(new File(filename));
    Set<Integer> seen = new HashSet<>();

    while (sc.hasNextLine()) {
        String line = sc.nextLine().trim();
        if (line.isEmpty()) continue;

        String[] parts = line.split("\\s+");
        if (parts.length < 4) continue;

        String edgeName  = parts[0];
        int    startNode = Integer.parseInt(parts[1]);
        int    endNode   = Integer.parseInt(parts[2]);
        double weight    = Double.parseDouble(parts[3]);

        graph.edges.add(new Edge(edgeName, startNode, endNode, weight));

        if (seen.add(startNode))
            graph.nodes.add(new Node(startNode, 0, 0));
        if (seen.add(endNode))
            graph.nodes.add(new Node(endNode,   0, 0));
    }
    return graph;
}
\end{lstlisting}

\subsection{Differences from the C implementation}

\begin{itemize}
  \item The C version uses a \textbf{two-pass} approach (validate first,
        allocate exact arrays, fill on second pass) because \texttt{malloc}
        requires a known size upfront. Java's \texttt{ArrayList} grows
        dynamically, so a single pass suffices.
  \item Error handling is done via \textbf{exceptions} (\texttt{FileNotFoundException},
        \texttt{NumberFormatException}) rather than the C \texttt{graph\_err\_t}
        enum with explicit exit codes.
  \item Duplicate-edge detection and connectivity checking (BFS) present in
        the C loader are not implemented in the Java version — the Swing GUI
        handles bad inputs gracefully via \texttt{JOptionPane} dialogs.
\end{itemize}

% ─────────────────────────────────────────────────────────────────────────────
\section{Layout Algorithms}

\subsection{Abstract Base — \texttt{LayoutAlgorithm}}

Both algorithms share a common interface declared in
\texttt{algorithm/LayoutAlgorithm.java}:

\begin{lstlisting}[caption={Abstract base class}]
public abstract class LayoutAlgorithm {
    public abstract void layout(Graph graph);
}
\end{lstlisting}

This mirrors the C approach of having a function pointer (or separate
function names) dispatched by the algorithm argument. In Java, the
\texttt{ToolPanel} holds a \texttt{currentAlgorithm} field of type
\texttt{LayoutAlgorithm} and polymorphically calls \texttt{layout()}.

\subsection{Fruchterman-Reingold (\texttt{Fr})}

The implementation in \texttt{algorithm/Fr.java} follows exactly the same
physical simulation model as the C version.

\subsubsection*{Constants}

\begin{center}
\begin{tabular}{lll}
\toprule
\textbf{Constant} & \textbf{Value} & \textbf{Purpose} \\
\midrule
\texttt{WIDTH}      & 800   & Canvas width (pixels) \\
\texttt{HEIGHT}     & 600   & Canvas height (pixels) \\
\texttt{PADDING}    & 30    & Margin from edge so labels are never clipped \\
\texttt{EPSILON}    & $10^{-4}$ & Guard against division by zero \\
\texttt{ITERATIONS} & 1000  & Fixed simulation length \\
\bottomrule
\end{tabular}
\end{center}

\subsubsection*{Step 1 — Initialisation}

A \texttt{HashMap<Integer,Integer>} maps each node's \texttt{id} to its
index in \texttt{graph.nodes}. This is needed because the edge list stores
IDs, not indices.

The optimal edge length and initial temperature are:
\[
  k = \sqrt{\frac{W \cdot H}{n}}, \qquad
  T_0 = \frac{W}{10}, \qquad
  \Delta T = \frac{T_0}{\text{ITERATIONS}}
\]

All nodes are placed at random positions inside the padded canvas using
\texttt{java.util.Random}.

\subsubsection*{Step 2 — Simulation loop (1000 iterations)}

\paragraph{A — Repulsive forces.}
Every pair $(i,j)$ is compared. The repulsive force magnitude is
\[
  f_r = \frac{k^2}{d_{ij}}
\]
and the displacement of node $i$ is incremented along the unit direction
$\hat{u}_{ij}$:

\begin{lstlisting}[caption={Repulsive force accumulation}]
double dx   = graph.nodes.get(i).X - graph.nodes.get(j).X;
double dy   = graph.nodes.get(i).Y - graph.nodes.get(j).Y;
double dist = Math.sqrt(dx*dx + dy*dy) + EPSILON;
double fr   = (k * k) / dist;
dispX[i] += (dx / dist) * fr;
dispY[i] += (dy / dist) * fr;
\end{lstlisting}

\paragraph{B — Attractive forces (weighted).}
Each edge pulls its endpoints together. The weight \texttt{e.len} scales
the force identically to the C version:
\[
  f_a = \frac{d^2}{k} \cdot w_e
\]

\begin{lstlisting}[caption={Attractive force accumulation}]
double fa    = (dist * dist / k) * e.len;
dispX[vi] -= (dx / dist) * fa;
dispY[vi] -= (dy / dist) * fa;
dispX[ui] += (dx / dist) * fa;
dispY[ui] += (dy / dist) * fa;
\end{lstlisting}

\paragraph{C — Temperature capping and boundary clamp.}
Displacement is limited to the current temperature, then the node is
clamped to the padded area:
\[
  \|\Delta \mathbf{p}_i\|_{\text{applied}} = \min(\|\Delta \mathbf{p}_i\|,\; T)
\]
\[
  x_i \leftarrow \text{clamp}(x_i,\; \text{PADDING},\; W - \text{PADDING})
\]

\paragraph{D — Cooling.}
At the end of each iteration: $T \leftarrow \max(0,\; T - \Delta T)$.

\textbf{Memory note.} Unlike C, no manual \texttt{free} is needed for the
displacement arrays \texttt{dispX}/\texttt{dispY}; they are local to
\texttt{layout()} and garbage-collected after the method returns.

\subsection{Tutte's Embedding (\texttt{Tutte})}

The implementation in \texttt{algorithm/Tutte.java} follows Tutte's
barycentric theorem.

\subsubsection*{Constants}

\begin{center}
\begin{tabular}{lll}
\toprule
\textbf{Constant} & \textbf{Value} & \textbf{Purpose} \\
\midrule
\texttt{RADIUS}    & $\min(W,H)\!\cdot\!0.4 - \text{PADDING}$ & Outer circle radius \\
\texttt{MAX\_ITER} & 1000  & Iteration limit \\
\texttt{TOLERANCE} & $10^{-6}$ & Convergence threshold \\
\bottomrule
\end{tabular}
\end{center}

\subsubsection*{Step 1 — Fix outer face}

The first $n_{\text{fixed}} = \min(3, n)$ nodes are pinned to a triangle
inscribed on a circle of radius \texttt{RADIUS} centred at $(W/2, H/2)$:
\[
  x_i = \frac{W}{2} + r\cos\!\left(\frac{2\pi i}{n_{\text{fixed}}}\right), \quad
  y_i = \frac{H}{2} + r\sin\!\left(\frac{2\pi i}{n_{\text{fixed}}}\right)
\]
All inner nodes start at the canvas centre $(W/2,\, H/2)$.

\subsubsection*{Step 2 — Adjacency list}

The Java version builds the adjacency structure using
\texttt{List<List<Integer>>} (an \texttt{ArrayList} of \texttt{ArrayList}s)
rather than the manually allocated 2-D C arrays. This eliminates the
\textit{two-pass degree count + malloc} pattern required in C.

\begin{lstlisting}[caption={Adjacency list construction}]
List<List<Integer>> adjList = new ArrayList<>();
for (int i = 0; i < n; i++) adjList.add(new ArrayList<>());

for (Edge e : graph.edges) {
    Integer a = idToIndex.get(e.startNode);
    Integer b = idToIndex.get(e.endNode);
    if (a == null || b == null) continue;
    adjList.get(a).add(b);
    adjList.get(b).add(a);
}
\end{lstlisting}

\subsubsection*{Step 3 — Gauss-Seidel relaxation}

Each inner node is moved to the arithmetic mean of its neighbours'
\emph{current} positions (in-place update, Gauss-Seidel style):
\[
  x_i \leftarrow \frac{1}{\deg(v_i)} \sum_{j \in N(i)} x_j, \quad
  y_i \leftarrow \frac{1}{\deg(v_i)} \sum_{j \in N(i)} y_j
\]

\subsubsection*{Step 4 — Convergence check}

After each node update, the squared displacement $\delta_i = (\Delta x_i)^2 +
(\Delta y_i)^2$ is tracked. The loop exits early when
$\max_i \delta_i < \varepsilon^2$.

\subsubsection*{Limitations}

Tutte's theorem guarantees a crossing-free embedding only for
\textbf{3-connected planar graphs}. For disconnected or non-planar graphs,
the output may contain edge crossings or overlapping nodes. The
implementation does not verify planarity or connectivity before running.

% ─────────────────────────────────────────────────────────────────────────────
\section{Graphical Interface (Swing)}

The GUI layer has no direct equivalent in the C project (which writes
coordinates to a file, then delegates visualisation to a separate Python
script). The Java version renders the graph in real time inside a
\texttt{JFrame} and supports interactive navigation.

\subsection{Window layout — \texttt{MainWindow}}

\texttt{MainWindow} extends \texttt{JFrame} and assembles the three main
panels using a \texttt{BorderLayout}:

\begin{center}
\begin{tabular}{lll}
\toprule
\textbf{Position} & \textbf{Component} & \textbf{Role} \\
\midrule
\texttt{CENTER} & \texttt{GraphPanel}      & Canvas — draws nodes and edges \\
\texttt{EAST}   & \texttt{ToolPanel}       & File loading, algorithm selection \\
\texttt{SOUTH}  & \texttt{CoordinatePanel} & Live cursor coordinates and zoom \\
\bottomrule
\end{tabular}
\end{center}

Window minimum size is $300 \times 200$\,px; default is $900 \times 600$\,px.

\subsection{Canvas — \texttt{GraphPanel}}

\texttt{GraphPanel} extends \texttt{JPanel} and overrides
\texttt{paintComponent} to render the graph. It maintains three
view-state variables:

\begin{center}
\begin{tabular}{lll}
\toprule
\textbf{Field} & \textbf{Type} & \textbf{Meaning} \\
\midrule
\texttt{scale}   & \texttt{double} & Current zoom factor (starts at \texttt{1.0}) \\
\texttt{offsetX} & \texttt{double} & Horizontal pan offset in screen pixels \\
\texttt{offsetY} & \texttt{double} & Vertical pan offset in screen pixels \\
\bottomrule
\end{tabular}
\end{center}

\subsubsection*{Coordinate transform}

Before drawing, an \texttt{AffineTransform} is applied so that graph
coordinates map to screen coordinates:
\[
  \mathbf{p}_{\text{screen}} = \begin{pmatrix} \texttt{offsetX} \\ \texttt{offsetY} \end{pmatrix}
  + \texttt{scale} \cdot \mathbf{p}_{\text{graph}}
\]

The inverse transform is used in mouse event handlers to convert a screen
click back to graph space:
\[
  g_x = \frac{e.x - \texttt{offsetX}}{\texttt{scale}}, \quad
  g_y = \frac{e.y - \texttt{offsetY}}{\texttt{scale}}
\]

\subsubsection*{Rendering pipeline}

\begin{enumerate}
  \item Apply \texttt{AffineTransform} to \texttt{Graphics2D} context.
  \item Draw all edges as \texttt{Line2D.Double} segments.
  \item Draw each node as a filled circle of radius \texttt{NODE\_RADIUS = 10}.
  \item Draw each node's ID as a string label offset by
        \texttt{(NODE\_RADIUS + 2, 4)} pixels.
\end{enumerate}

\subsection{Interaction}

\subsubsection*{Zoom}

Mouse wheel events call \texttt{onScroll}. The zoom factor per scroll
tick is \texttt{ZOOM\_FACTOR = 1.03} (3\%). To keep the point under the
cursor stationary during zoom, the offset is adjusted:

\begin{lstlisting}[caption={Zoom centred on cursor}]
double factor = (e.getPreciseWheelRotation() < 0) ? ZOOM_FACTOR : 1.0/ZOOM_FACTOR;
offsetX = e.getX() - factor * (e.getX() - offsetX);
offsetY = e.getY() - factor * (e.getY() - offsetY);
scale  *= factor;
\end{lstlisting}

\subsubsection*{Pan}

Dragging on an empty area (no node selected) translates the canvas by
updating \texttt{offsetX} and \texttt{offsetY}:
\[
  \texttt{offsetX} \mathrel{+}= \Delta x_{\text{mouse}}, \quad
  \texttt{offsetY} \mathrel{+}= \Delta y_{\text{mouse}}
\]

\subsubsection*{Node dragging}

On \texttt{mousePressed}, the code finds the closest node to the cursor
in graph space. If the distance is less than \texttt{NODE\_RADIUS / scale}
(hit-test scaled to current zoom), that node is stored in
\texttt{selectedNode}. Subsequent drag events move the node directly:
\[
  \texttt{node.X} \mathrel{+}= \frac{\Delta x_{\text{mouse}}}{\texttt{scale}}, \quad
  \texttt{node.Y} \mathrel{+}= \frac{\Delta y_{\text{mouse}}}{\texttt{scale}}
\]

\subsection{Tool Panel — \texttt{ToolPanel}}

\texttt{ToolPanel} (right side, \texttt{150 px} wide) provides:

\begin{itemize}
  \item \textbf{Select File} — opens a \texttt{JFileChooser} starting from
        \texttt{./test-graphs/} (or the last opened directory). On
        successful load, the chosen algorithm's \texttt{layout()} is run
        immediately.
  \item \textbf{Algorithm radio buttons} — \textit{Tutte} (default) and
        \textit{FR}. Switching re-runs the layout on the current graph.
  \item \textbf{Export PNG} — stub, not yet implemented.
  \item \textbf{Reset View} — calls \texttt{GraphPanel.resetView()},
        restoring \texttt{scale=1}, \texttt{offsetX=offsetY=0}.
\end{itemize}

\textbf{Design note.} The \texttt{tutteRadio.setSelected(true)} call is
made \emph{before} the \texttt{ItemListener} is attached. This is
intentional: attaching the listener first would trigger a layout run
before any graph has been loaded.

\subsection{Status Bar — \texttt{CoordinatePanel}}

A thin strip at the bottom of the window displays three live values:

\begin{itemize}
  \item \textbf{Coordinate X / Y} — graph-space position of the mouse cursor,
        updated on every \texttt{mouseMoved} event.
  \item \textbf{Zoom} — current \texttt{scale} factor, updated on every
        scroll event.
\end{itemize}

% ─────────────────────────────────────────────────────────────────────────────
\section{Input File Format}

The file format is identical to the C version: a plain-text edge list where
each non-empty line describes one undirected edge as a 4-tuple:

\[
  \langle \textit{Name},\; \textit{ID}_A,\; \textit{ID}_B,\; \textit{Weight} \rangle
\]

\begin{center}
\begin{tabular}{llll}
\toprule
\textbf{Field}  & \textbf{Type}   & \textbf{Constraints} & \textbf{Example} \\
\midrule
Name     & String  & No whitespace          & \texttt{L0\_12} \\
$ID_A$   & int     & Positive integer       & \texttt{1} \\
$ID_B$   & int     & Positive integer       & \texttt{2} \\
Weight   & double  & Any double             & \texttt{1.0} \\
\bottomrule
\end{tabular}
\end{center}

Fields are separated by whitespace. Blank lines and lines with fewer than
4 fields are silently skipped by \texttt{GraphReader}. The vertex set is
inferred automatically from the IDs encountered in the edge list — no
explicit vertex count is required.

\subsubsection*{Example}

\begin{lstlisting}[style=plain,caption={test.txt --- a Petersen-like graph}]
L0_12   1  2  1.0
L0_23   2  3  1.0
L0_13   1  3  1.0
E1_14   1  4  1.0
E1_25   2  5  1.0
...
C_910   9  10 1.0
\end{lstlisting}

This file defines 10 vertices and 24 edges. Vertex IDs need not be
contiguous or start from 1; the reader handles arbitrary positive integers.

\subsubsection*{Output}

\texttt{GraphWriter} is currently a stub class. Coordinate output (the
equivalent of the C \texttt{output\_write} function) has not been
implemented; the primary output medium is the interactive Swing canvas.
The \textbf{Export PNG} button in \texttt{ToolPanel} is also a planned
but unfinished feature.

% ─────────────────────────────────────────────────────────────────────────────
\section{Build and Run}

\subsection{Requirements}

\begin{itemize}
  \item Java 17 or newer (any distribution: OpenJDK, Temurin, etc.)
  \item An IDE such as IntelliJ IDEA is recommended; the \texttt{P2.iml}
        module file is already included.
\end{itemize}

\subsection{Running from the command line}

\begin{lstlisting}[style=bash,caption={Compile and run from project root}]
# Compile all sources into out/
javac -d out/ src/Main.java \
      src/model/*.java      \
      src/io/*.java         \
      src/algorithm/*.java  \
      src/gui/*.java

# Run
java -cp out Main
\end{lstlisting}

The application opens a $900 \times 600$\,px window. Use
\textbf{Select File} in the right panel to load any \texttt{.txt}
edge-list file, then choose a layout algorithm.

\subsection{Running from IntelliJ IDEA}

\begin{enumerate}
  \item Open the project root directory (\texttt{File → Open}).
  \item Mark \texttt{src/} as the Sources Root
        (right-click → \textit{Mark Directory as → Sources Root}).
  \item Run the \texttt{Main} class directly (\texttt{Shift+F10}).
\end{enumerate}

% ─────────────────────────────────────────────────────────────────────────────
\section{Java vs.\ C — Design Comparison}

\begin{center}
\begin{tabular}{p{3.5cm}p{5cm}p{5cm}}
\toprule
\textbf{Aspect} & \textbf{C implementation} & \textbf{Java implementation} \\
\midrule
Data structures  &
  Structs (\texttt{vertex\_t}, \texttt{edge\_t}, \texttt{graph\_t}) with
  explicit size counters &
  Classes (\texttt{Node}, \texttt{Edge}, \texttt{Graph}) with
  \texttt{ArrayList} \\
\addlinespace
Memory management &
  Manual: \texttt{malloc}/\texttt{calloc}/\texttt{free};
  two-pass load to compute exact sizes &
  Automatic (GC); single-pass load with dynamic lists \\
\addlinespace
Error handling &
  \texttt{graph\_err\_t} enum + process exit codes &
  Java exceptions + \texttt{JOptionPane} dialogs \\
\addlinespace
Algorithm dispatch &
  Conditional branch on a string argument &
  Polymorphism via \texttt{LayoutAlgorithm} subclasses \\
\addlinespace
Visualisation &
  Coordinates written to file; Python + matplotlib renders PNG &
  Real-time interactive Swing canvas \\
\addlinespace
Adjacency list &
  Two-pass: count degrees, \texttt{malloc}, fill &
  \texttt{ArrayList<ArrayList<Integer>>} built in one pass \\
\bottomrule
\end{tabular}
\end{center}

% ─────────────────────────────────────────────────────────────────────────────
\end{document}
